\section{Lean formalization architecture}
\label{sec:lean-plan}

This section gives a source-oriented plan against the checked \texttt{ConciseIPRA} API.  It is deliberately more precise than a list of theorem names: it fixes variable order, formula shapes, module dependencies, and the point at which each semantic construction becomes an actual object-iPRA derivation.

\subsection{Pinned upstream interface}

The intended imports are the repository modules providing:

\begin{itemize}
\item \texttt{Syn}, \texttt{Fm}, substitution, shifting, and de Bruijn utilities;
\item \texttt{Prf} and notation \texttt{Gamma |-[PRAAx] phi};
\item primitive-recursive syntax \texttt{PR} and terms \texttt{PTm};
\item the iPRA axiom family \texttt{PRAAx} and \texttt{PRA.qfInd};
\item the native transfer theorem \texttt{pra\_iff\_native}, used only after an ordinary derivation is complete.
\end{itemize}

The inspected repository examples import \path{ConciseIPRA.PR} and work in namespace \path{LeanProofs.ConciseIPRA}.  The packet's blueprint follows that convention.  Before integration, the exact commit should be pinned and every imported name checked against the current module tree.

\subsection{Module graph}

The proposed development is split into ten modules.

\begin{longtable}{@{}p{0.27\textwidth}p{0.65\textwidth}@{}}
\toprule
Module & Responsibility \\
\midrule
\endhead
\texttt{PFIG/Boolean.lean} & Boolean terms, range theorems, reflection of connectives, bounded \texttt{all}/\texttt{any}, least search, and QF formula Booleanization. \\
\path{PFIG/ExactContract.lean} & Exact term contracts, specialization, composition, and stable-observer combinators.  Most proofs adapt existing repository code. \\
\texttt{PFIG/ObjectLoop.lean} & Source-level wrapper around \texttt{primitiveRecursionInvariant} and \texttt{primitiveRecursionAt}; terminal extraction. \\
\texttt{PFIG/Demand.lean} & Coded request semilattices, monotone demand maps, fixed-DAG reverse pass, and object reflection of demand equations. \\
\texttt{PFIG/Relation.lean} & Finite relation codes, viability table layout, path checker, least viable path, and object reflection. \\
\texttt{PFIG/Shrinking.lean} & Coded rational metrics, viable-cloud diameter, edge motion, adjacent coherence, finite triangle sums, and regular-name extraction. \\
\texttt{PFIG/Guarded.lean} & Apartness guards, totalized operations, branch stabilization, observational setoids, and relational back-and-forth. \\
\texttt{PFIG/Graph.lean} & Fixed graph syntax or generated-stage syntax, topological schedule, two-pass state compiler, invariant assembly, and soundness theorem. \\
\path{PFIG/NegativeTransfer.lean} & Equality stability, stability of negative translations, proof translation, retained equivalence, and retained transfer. \\
\texttt{PFIG/Applications.lean} & Worked end-to-end examples and adapters for height/list, finite gluing, root clusters, and finite nets. \\
\bottomrule
\end{longtable}

The first complete milestone need not internalize the graph AST.  It may use generated theorem-specific stages and prove them with \texttt{ObjectLoop}.  \texttt{Graph.lean} can then abstract the repeated pattern.

\subsection{Exact contracts}

The exact interface should reuse the repository's type:

\begin{lstlisting}[language={},caption={Exact object contract}]
structure ExactContract {n : Nat}
    (P : Fm PTm n) (Q : Fm PTm (n + 1)) where
  witness   : PTm n
  certified : [P] |-[PRAAx] Q.subst (inst witness)
\end{lstlisting}

Required derived operations are:

\begin{lstlisting}[language={}]
theorem ExactContract.existsOutput :
  c : ExactContract P Q -> [P] |-[PRAAx] Fm.ex Q

def ExactContract.compose
    (c : ExactContract P Q)
    (d : ExactContract Q R) :
    ExactContract P (R.subst (up (inst c.witness)))

def ExactContract.specialize
    (c : ExactContract P Q) (sigma : Fin n -> PTm m) :
    ExactContract (P.subst sigma) (Q.subst (up sigma))
\end{lstlisting}

The inspected \texttt{ObjectComposition.lean} already gives these constructions.  The new module should import or relocate them rather than duplicate proofs.

\subsection{Loop contracts}

Use a formula \(P\) with free variables ordered as state, time, then parameters.  Define helper substitutions exactly once.

\begin{lstlisting}[language={}]
def instStateTime {k n : Nat}
    (P : Fm PTm (k + 2))
    (state time : PTm n)
    (params : Fin k -> PTm n) : Fm PTm n :=
  P.subst (Fin.cases state (Fin.cases time params))
\end{lstlisting}

A source-facing contract is:

\begin{lstlisting}[language={}]
structure LoopContract {k : Nat}
    (Gamma : List (Fm PTm k)) (Post : Fm PTm k) where
  init      : PR k
  step      : PR (k + 2)
  fuel      : PTm k
  invariant : Fm PTm (k + 2)
  qf        : invariant.QF
  base      : Gamma |-[PRAAx] baseFormula init invariant
  preserve  : invariant :: (Gamma.map Fm.shift).map Fm.shift
              |-[PRAAx] nextFormula step invariant
  finish    : Gamma |-[PRAAx]
              instStateTime invariant
                (runTerm init step fuel var) fuel var --> Post
\end{lstlisting}

Then:

\begin{lstlisting}[language={}]
theorem LoopContract.sound (c : LoopContract Gamma Post) :
  Gamma |-[PRAAx] Post :=
  Prf.impE c.finish
    (primitiveRecursionAt c.init c.step c.invariant
      c.qf Gamma c.base c.preserve c.fuel)
\end{lstlisting}

All substitution lemmas required by \texttt{primitiveRecursionAt} already appear in the inspected \texttt{ObjectIteration.lean}.  The integration task is to make that theorem a stable library module and remove copied \texttt{ObjectKernel} snapshots from individual packets.

\subsection{Boolean formulas}

There are two useful layers.

\paragraph{Term-level Boolean algebra.}
Fix PR symbols \texttt{bnot}, \texttt{band}, \texttt{bor}, \texttt{beq}, and bounded folds.  Prove object theorems such as:

\begin{lstlisting}[language={}]
theorem band_reflect (a b : PTm n) :
  [] |-[PRAAx]
    ((band a b = 1) <-> ((a = 1) /\ (b = 1)))

theorem bool_range (a : PTm n) :
  BoolTerm a -> [] |-[PRAAx] ((a = 0) \/ (a = 1))
\end{lstlisting}

The concrete formula constructors use the repository's notation, not the ASCII shown in the schematic listing.

\paragraph{Syntax-directed QF compiler.}
Define a metalevel function from a quantifier-free \(\Fm\) to a PR Boolean expression.  Because arbitrary formula syntax contains term symbols already interpreted by \texttt{PTm}, the compiler may return a \texttt{PTm n} plus a derivation record:

\begin{lstlisting}[language={}]
structure BoolReflection {n : Nat} (A : Fm PTm n) where
  code    : PTm n
  range   : [] |-[PRAAx] (code = 0 \/ code = 1)
  reflect : [] |-[PRAAx] Fm.iff (code = 1) A

def booleanizeQF {n : Nat} (A : Fm PTm n) (h : A.QF) :
  BoolReflection A
\end{lstlisting}

For a fully generic compiler, implication and disjunction cases use the already constructed decidability of subformulas.  A simpler first milestone provides reflection constructors and lets generated files assemble them.

\subsection{Bounded search}

Define a fixed PR state encoding a found bit and witness.  The theorem should expose both correctness and minimality:

\begin{lstlisting}[language={}]
structure SearchSpec {n : Nat} where
  bound : PTm n
  pred  : Fm PTm (n + 1)
  qf    : pred.QF
  bool  : BoolReflection pred

def leastSearch (s : SearchSpec) : PTm n

theorem leastSearch_sound
    (nonempty : Gamma |-[PRAAx] Fm.exBounded s.bound s.pred) :
  Gamma |-[PRAAx]
    s.pred.subst (inst (leastSearch s))

theorem leastSearch_minimal ...
\end{lstlisting}

If \texttt{Fm.exBounded} is not a library macro, define it as an ordinary existential conjoined with \(<\), and separately use the Booleanized bounded form for the loop invariant.

\subsection{Finite relations and table layout}

Keep object formulas independent of semantic finite types.  Use natural bounds and a fixed table encoding.

\begin{lstlisting}[language={}]
structure RelPresentation (n : Nat) where
  candBound : PTm (n + 1)       -- precision, params
  good      : Fm PTm (n + 2)    -- candidate, precision, params
  edge      : Fm PTm (n + 3)    -- next, current, precision, params
  goodQF    : good.QF
  edgeQF    : edge.QF
\end{lstlisting}

Document the variable order in comments and provide named substitution functions:

\begin{lstlisting}[language={}]
def RelPresentation.goodAt
    (R : RelPresentation n)
    (a p : PTm m) (xs : Fin n -> PTm m) : Fm PTm m

def RelPresentation.edgeAt
    (R : RelPresentation n)
    (b a p : PTm m) (xs : Fin n -> PTm m) : Fm PTm m
\end{lstlisting}

The viability table stores one bit per pair \((i,a)\) with \(i\le H\) and \(a<B_i\).  To avoid variable-length arrays in the object proof, use a rectangular bound
\[
 B^*(H,\vec{x})=\max_{i\le H}B_i(\vec{x})
\]
and mark out-of-range cells false.  Index \((i,a)\) by a primitive-recursive pairing function.  Reverse iteration can be implemented by the index \(H-i\).

Target theorems are:

\begin{lstlisting}[language={}]
theorem viableTable_reflect :
  Gamma |-[PRAAx]
    TableOK R H table -->
    (tableBit table i a = 1 <-> ViableFormula R H i a)

theorem leastViablePath_sound :
  Gamma |-[PRAAx] NonemptyViable R H -->
    PathOK R H (leastViablePath R H)
\end{lstlisting}

\texttt{ViableFormula} may contain bounded existentials and is used only after Boolean reflection; \texttt{TableOK} and the induction invariant remain QF.

\subsection{Shrinking metrics}

Represent nonnegative rationals by numerator/positive-denominator pairs or scaled integers.  Define a metric observation interface:

\begin{lstlisting}[language={}]
structure ObservedRelation (n : Nat) extends RelPresentation n where
  obs       : PR (n + 2)    -- candidate, precision, params -> rational code
  distLe    : Fm PTm (n + 3)
  diameter : PTm (n + 1)    -- precision, params -> rational code
  motion   : PTm (n + 1)
  diameterProof : ...       -- actual Prf PRAAx derivation
  motionProof   : ...
\end{lstlisting}

The exact arity of \texttt{distLe} depends on the rational-code API; the blueprint file keeps the theorem signature schematic.  The key object targets are:

\begin{lstlisting}[language={}]
theorem adjacent_coherence :
  Gamma |-[PRAAx]
    distLe (x p) (x (p+1)) (diameter p + motion p)

theorem regularName_of_tail :
  Gamma |-[PRAAx] TailBudget diameter motion --> RegularName x
\end{lstlisting}

The finite triangle-sum theorem is proved by \texttt{PRA.qfInd} on an equality/inequality code, not by invoking a semantic metric-space hierarchy.

\subsection{Observational setoids}

Use a structure containing formulas rather than Lean quotient types:

\begin{lstlisting}[language={}]
structure Presentation (n : Nat) where
  valid  : Fm PTm (n + 1)          -- code, params
  equiv  : Fm PTm (n + 3)          -- x, y, precision, params
  validQF : valid.QF
  equivQF : equiv.QF
  refl   : ...
  symm   : ...
  trans  : ...
  restrict : PR ...
  restrict_sound : ...
\end{lstlisting}

An operation record stores a code function, demand, guard, correctness, and extensionality derivation.  Relational operations replace a code function by a finite relation and extensionality by bounded back-and-forth.

\subsection{Demand semilattices}

Start with scalar requests \(\Nat\), where join is maximum.  After the scalar compiler works, generalize to coded tuples:

\begin{lstlisting}[language={}]
structure RequestCode where
  join : PR 2
  le   : PR 2
  refl : ...
  join_left  : ...
  join_right : ...
  join_least : ...
\end{lstlisting}

A fixed graph can initially carry plain Lean functions producing \texttt{PTm} demand expressions.  The object reverse pass is generated as one loop with a QF invariant.  Internalizing arbitrary graph syntax is not required for theorem-level soundness.

\subsection{Graph syntax and compilation}

After generated examples stabilize, define a typed metalevel graph.  One workable design separates numerical environments from proof obligations:

\begin{lstlisting}[language={}]
inductive NodeKind
  | exact | loop | search | relation | guard | observe | glue

structure Node where
  id        : Nat
  kind      : NodeKind
  deps      : List Nat
  request   : RequestSchema
  stateSize : Nat
  localCode : LocalProgram
  contract  : LocalContractSchema

structure Graph where
  nodes : List Node
  acyclic : AcyclicByIndex nodes
\end{lstlisting}

\texttt{LocalProgram} is not an arbitrary PR code evaluated internally.  It is metasyntax elaborated to a concrete \texttt{PR} term.  The compiler returns a package:

\begin{lstlisting}[language={}]
structure CompiledGraph (G : Graph) where
  demandInit demandStep demandFuel : ...
  execInit   execStep   execFuel   : ...
  demandInv execInv : Fm PTm ...
  demandQF : demandInv.QF
  execQF   : execInv.QF
  demandBase demandPreserve : Prf PRAAx ...
  execBase execPreserve     : Prf PRAAx ...
  extract : Prf PRAAx ...
\end{lstlisting}

\begin{lstlisting}[language={}]
theorem CompiledGraph.sound (C : CompiledGraph G) :
  InputContext G |-[PRAAx] TerminalFormula G := ...
\end{lstlisting}

The proof invokes the two loop soundness theorems and cuts \texttt{extract}.  Structural soundness of the compiler is an induction on \texttt{G.nodes}; acyclicity by index makes dependency lookup simple.

\subsection{Negative translation modules}

Use the existing definitions from \texttt{NegativeSyntax.lean}.  The target signatures are:

\begin{lstlisting}[language={}]
theorem equality_stable (s t : PTm n) :
  [] |-[PRAAx] ~~(s = t) --> (s = t)

theorem negative_stable (A : Fm PTm n) :
  [] |-[PRAAx] ~~(negative A) --> negative A

theorem translate_prf
    (h : Gamma |-[ClassicalPRAAx] A) :
  Gamma.map negative |-[PRAAx] negative A

theorem retained_iff
    (h : Retained A) :
  [] |-[PRAAx] Fm.iff (negative A) A

theorem retained_transfer
    (hGamma : forall B in Gamma, Retained B)
    (hA : Retained A) :
  (Gamma |-[ClassicalPRAAx] A) <-> (Gamma |-[PRAAx] A)
\end{lstlisting}

The actual syntax for double negation is \texttt{negFormula (negFormula A)}.  The \texttt{orE} and \texttt{exE} branches of \texttt{translate\_prf} should be implemented with helper lemmas:

\begin{lstlisting}[language={}]
theorem neg_or_elim
    (hOr : Gamma |- negative (A or B))
    (hA  : negative A :: Gamma |- C)
    (hB  : negative B :: Gamma |- C)
    (stableC : Gamma |- ~~C --> C) : Gamma |- C

theorem neg_ex_elim ...
\end{lstlisting}

This isolates the binder and context bookkeeping from the main induction.

\subsection{Worked formalization milestones}

The recommended order is proof-driven.

\begin{enumerate}
\item Re-export the existing exact contract, object iteration, and stable-observer files in one library without copied kernels.
\item Complete Boolean range/reflection for a small fixed set of PR Boolean symbols.
\item Formalize least bounded search and one example: choose the least divisor or least valid rational interval.
\item Formalize a tiny viability table with fixed level bounds, then generalize to variable bounds.
\item Prove shrinking coherence for exact rational observations and a geometric tail budget.
\item Generate and check one two-pass PFIG: bisection for \(y^3+y=x\), including backward precision demand and a regular root name.
\item Complete negative translation and retained transfer.
\item Add finite gluing and height/list adapters.
\item Only then internalize the general graph AST and prove compiler soundness.
\end{enumerate}

This order yields useful checked theorems early and keeps the generic compiler grounded in concrete examples.

\subsection{Axiom and status audit}

Every public theorem concluding \(\Prf\ \PRAAx\) should be followed in an audit file by \texttt{\#print axioms}.  Expected metatheoretic axioms inherited from Lean's quotient/propositional infrastructure should be recorded, and no new arithmetic axiom should be introduced.  Files containing \texttt{sorry}, \texttt{axiom}, or theorem propositions used as placeholders must be clearly marked blueprint-only.

The build should produce a machine-readable status file with, for each theorem:

\begin{itemize}
\item semantic Lean checked;
\item executable tests passed;
\item object formula defined;
\item \(\Prf\ \PRAAx\) derivation checked;
\item native transfer checked;
\item remaining producer assumptions.
\end{itemize}

This prevents a semantic theorem from being mistaken for completion of the object layer.

\subsection{Clear path to a complete formalization}

No new foundational principle is required.  The completion path consists of:

\begin{enumerate}
\item consolidating already existing object-level loop and contract code;
\item implementing finite Boolean reflection and table libraries;
\item proving viability and shrinking estimates by those libraries;
\item generating fixed PR state transitions and QF invariants;
\item composing local object certificates by substitution and cut;
\item completing the standard negative-translation induction;
\item applying the already checked ordinary-to-native equivalence at the end.
\end{enumerate}

The largest engineering risk is substitution bookkeeping for generated formulas, not proof-theoretic strength.  Named instantiation helpers, a fixed variable-order convention, and small regression examples should make that risk manageable.
